\documentclass[11pt,a4paper]{article}

\usepackage[margin=2.5cm]{geometry}
\usepackage{booktabs}
\usepackage{graphicx}
\usepackage{hyperref}
\usepackage{amsmath}
\usepackage{array}
\usepackage{xcolor}
\usepackage{enumitem}
\usepackage{caption}

\hypersetup{colorlinks=true, linkcolor=blue!60!black, urlcolor=blue!60!black}

\title{%
  \textbf{Progress Report: Two-Stream Action Recognition on Charades}\\[0.4em]
  \large CNN Baselines and VideoMAE\,V2 Fine-Tuning (Fir \& Nibi)
}
\author{Asma\\
  \textit{Supervisor: Prof.\ \'Eric Hervet}\\
  \textit{Alliance Canada --- Mitacs project}}
\date{July 7, 2026}

\begin{document}
\maketitle

%==============================================================================
\section{Experimental Setup}
%==============================================================================

\subsection{Dataset and evaluation protocol}

\textbf{Charades} multi-label action recognition (157 classes). Inputs: RGB frames
($224\times224$) and precomputed AllenAI optical flow ($x/y$ JPEG pairs). Labels from
\texttt{Charades\_v1\_\{train,test\}.csv}; each video may have multiple simultaneous
actions. Compact subsets are staged from \texttt{\$SCRATCH} to node-local SSD per job;
frame-count sweeps use denser 6\,fps or 12\,fps archives
(\texttt{charades\_2stream\_subset\_\{6,12\}fps.tar}).

All reported runs use the \textbf{\texttt{train\_test}} protocol: train on the official
training set; select checkpoints and report final metrics on the official \textbf{test}
set. 

\textbf{Metrics:} macro mAP (unweighted mean of per-class AP; Charades standard) and
weighted mAP (AP weighted by positive count). \textbf{Loss:} BCE with sigmoid outputs;
CNN training optionally uses per-class positive weighting; VideoMAE uses unweighted BCE.

\subsection{Model architectures}

\begin{table}[h]
\centering
\caption{CNN two-stream architecture (\texttt{src/models.py}, Fir baselines).}
\begin{tabular}{@{}ll@{}}
\toprule
Component & Details \\
\midrule
RGB stream & ResNet-50 (ImageNet-1K), temporal mean pooling over frames \\
Flow stream & ResNet-50, \texttt{conv1} inflated to $2 \times \texttt{num\_flow}$ channels \\
Fusion & Late fusion: \texttt{average}, \texttt{weighted}, or \texttt{concat} \\
Projection & Linear $\rightarrow$ 512-d ReLU before classifier \\
Classifier & 157 logits (sigmoid at inference) \\
Temporal chunking & 16 frames per ResNet forward pass (VRAM control) \\
Grad checkpointing & Enabled for dense/long-frame sweeps \\
\bottomrule
\end{tabular}
\end{table}

\begin{table}[h]
\centering
\caption{VideoMAE two-stream architecture (\texttt{src/videomae\_model.py}, Nibi).}
\begin{tabular}{@{}ll@{}}
\toprule
Component & Details \\
\midrule
RGB backbone & VideoMAE ViT-Base (\texttt{vit\_base\_patch16\_224}), 3 input channels \\
Flow backbone & Same ViT-Base architecture, 2 input channels \\
Pretraining & \texttt{vit\_b\_k710\_dl\_from\_giant.pth} (Kinetics-710, VideoMAEv2) \\
Flow init & Flow backbone initialized from RGB checkpoint (first 2 patch-embed channels) \\
Fusion & \texttt{concat} $\rightarrow$ MLP (1536 $\rightarrow$ 512 $\rightarrow$ 157) \\
Tubelet size & 2; image size $224 \times 224$ \\
\bottomrule
\end{tabular}
\end{table}

Flow \texttt{patch\_embed} is frozen for the first $N$ epochs (3 in the frame sweep,
5 in the reproducible recipe), then unfrozen.

\subsection{Compute infrastructure}

\begin{table}[h]
\centering
\caption{Cluster configuration.}
\begin{tabular}{@{}lll@{}}
\toprule
Cluster & GPU & Primary role in this report \\
\midrule
\textbf{Fir} & NVIDIA H100 (80\,GB) & CNN baselines + VideoMAE RGB fine-tuning \\
\textbf{Nibi} & NVIDIA H100 (80\,GB) & VideoMAE two-stream fine-tuning (in progress) \\
\bottomrule
\end{tabular}
\end{table}

Both clusters use Python~3.11, PyTorch, and \texttt{charades\_venv}. Data and pretrained
checkpoints reside on \texttt{\$SCRATCH/charades/}.

\subsection{Hyperparameters --- CNN baselines (Fir)}

Table~\ref{tab:cnn-hparams} lists hyperparameters across CNN experiment groups. Default
RGB/flow/fusion jobs use the standard 16-frame subset
(\texttt{charades\_2stream\_subset.tar}, $\sim$16 RGB frames and 10 flow pairs per video).

\begin{table}[h]
\centering
\caption{CNN training hyperparameters by experiment group.}
\label{tab:cnn-hparams}
\small
\begin{tabular}{@{}lcccccc@{}}
\toprule
Experiment group & Epochs & Batch & LR & Weight decay & RGB frames & Flow pairs \\
\midrule
\texttt{rgb}, \texttt{flow}, \texttt{fusion\_*} (default) & 25 & 32 / 24 & $10^{-4}$ & $10^{-4}$ & 8 & 10 \\
\texttt{rgb\_n\{2,4,8,12,16\}\_ctrl} & 25 & 32 & $10^{-4}$ & $10^{-4}$ & 2--16 & 10 \\
\texttt{rgb\_n\{24,48,96\}\_dense\_e20} & 20 & 16 & $10^{-4}$ & $10^{-4}$ & 24--96 & 10 \\
\texttt{rgb\_n\{32,64\}\_dense\_e20} & 20 & 16 / 8 & $10^{-4}$ & $10^{-4}$ & 32 / 64 & 10 \\
\texttt{flow\_n\{32,64\}\_dense\_e20} & 20 & 16 / 8 & $10^{-4}$ & $10^{-4}$ & 8 & 32 / 64 \\
\texttt{fusion\_*\_n\{32,64\}\_dense\_e20} & 20 & 6 / 3 & $10^{-4}$ & $10^{-4}$ & 32 / 64 & 32 / 64 \\
\bottomrule
\end{tabular}

\vspace{0.5em}
\footnotesize
All dense sweeps: 12\,fps subset, \texttt{temporal\_chunk=16}, gradient checkpointing enabled.
Fusion variants: \texttt{average}, \texttt{weighted}, \texttt{concat}.
\texttt{proj\_dim=512}, \texttt{img\_size=224}, BCE with positive class weighting.
\end{table}

\subsection{Hyperparameters --- VideoMAE RGB (Fir)}

Completed RGB-only fine-tuning via \texttt{jobs/fir/train\_videomae.sh} on the standard
16-frame subset (\texttt{charades\_2stream\_subset.tar}).

\begin{table}[h]
\centering
\caption{VideoMAE RGB fine-tuning hyperparameters (Fir, \texttt{videomae\_vitb}).}
\label{tab:vmae-rgb-hparams}
\small
\begin{tabular}{@{}ll@{}}
\toprule
Parameter & Value \\
\midrule
Model & \texttt{vit\_base\_patch16\_224} (3 input channels) \\
Pretrained checkpoint & \texttt{vit\_b\_k710\_dl\_from\_giant.pth} \\
Frames per clip & 16 (all stored subset frames) \\
Epochs & 30 \\
Batch size & 16 \\
Peak learning rate & $10^{-3}$ (layer-wise decay 0.75) \\
Warmup epochs & 2 (cosine to $10^{-6}$) \\
Weight decay / drop path & 0.05 / 0.1 \\
Loss & BCE (unweighted) \\
Protocol & \texttt{train\_test} \\
\bottomrule
\end{tabular}
\end{table}

\subsection{Hyperparameters --- VideoMAE two-stream (Nibi)}

The Nibi frame sweep was launched via \texttt{jobs/nibi/sweep\_videomae\_twostream\_frames.sh}
(array tasks for 16 / 32 / 48 / 64 frames). Table~\ref{tab:vmae-hparams} contrasts the
\textbf{sweep recipe} (actually used) with the \textbf{reproducible recipe}
(\texttt{jobs/nibi/train\_videomae\_twostream.sh}).

\begin{table}[h]
\centering
\caption{VideoMAE two-stream hyperparameters.}
\label{tab:vmae-hparams}
\small
\begin{tabular}{@{}lcc@{}}
\toprule
Parameter & Frame sweep (Nibi, used) & Reproducible recipe (f16) \\
\midrule
Model & \texttt{vit\_base\_patch16\_224} $\times$ 2 & same \\
Pretrained checkpoint & \texttt{vit\_b\_k710\_dl\_from\_giant.pth} & same \\
Fusion & \texttt{concat} & \texttt{concat} \\
Epochs (target) & 30 & 30 \\
Warmup epochs & 2 & 2 \\
Peak learning rate & $\mathbf{10^{-3}}$ & $\mathbf{5 \times 10^{-4}}$ \\
Min learning rate & $10^{-6}$ (cosine) & $10^{-6}$ \\
Layer-wise LR decay & 0.75 & 0.75 \\
Weight decay & 0.05 & 0.05 \\
Drop path & 0.1 & 0.1 \\
Gradient clipping & --- & 1.0 \\
Freeze flow embed & 3 epochs & 5 epochs \\
Random seed & 42 (default) & 42 \\
Loss & BCE (unweighted) & BCE (unweighted) \\
Protocol & \texttt{train\_test} & \texttt{train\_test} \\
Subset & 6\,fps dense (sweep) & 16-frame standard subset \\
\midrule
Batch size (16 fr.) & 6 & 6 \\
Batch size (32 fr.) & 4 (script) / \textbf{2 (actual)} & 4 \\
Batch size (64 fr.) & 2 (script) / \textbf{1 (actual)} & 2 \\
\bottomrule
\end{tabular}
\end{table}

\noindent
\textbf{Note:} Actual batch sizes for f32 and f64 (2 and 1) differ from script defaults,
indicating manual reduction after out-of-memory (OOM) errors.

%==============================================================================
\section{Results}
%==============================================================================

\noindent
Table~\ref{tab:headline} highlights the strongest completed runs at comparable settings;
full CNN sweeps follow in Table~\ref{tab:fir-results}.

\begin{table}[h]
\centering
\caption{Headline results }
\label{tab:headline}
\small
\begin{tabular}{@{}llrr@{}}
\toprule
Run & Cluster & Macro mAP & Weighted mAP \\
\midrule
\textbf{\texttt{videomae\_vitb}} (16 RGB fr., 30 ep.) & Fir & \textbf{0.3259} & \textbf{0.4238} \\
\texttt{rgb\_n32\_dense\_e20} (32 RGB fr., 20 ep.) & Fir & 0.1701 & 0.2380 \\
\texttt{rgb\_n16\_ctrl} (16 RGB fr., 25 ep.) & Fir & 0.1690 & 0.2338 \\
\texttt{fusion\_weighted} (8 RGB + 10 flow) & Fir & 0.1649 & 0.2307 \\
\texttt{flow} (10 flow pairs) & Fir & 0.0989 & 0.1544 \\
\bottomrule
\end{tabular}
\end{table}

\subsection{CNN baselines on Fir (completed runs)}

\begin{table}[h]
\centering
\caption{CNN baseline results on Fir }
\label{tab:fir-results}
\small
\begin{tabular}{@{}lrr@{}}
\toprule
Run & Macro mAP & Weighted mAP \\
\midrule
\multicolumn{3}{@{}l}{\textit{Single-stream defaults (16-frame subset)}} \\
\texttt{rgb} & 0.1616 & 0.2243 \\
\texttt{flow} & 0.0989 & 0.1544 \\
\midrule
\multicolumn{3}{@{}l}{\textit{Fusion (default frames)}} \\
\texttt{fusion\_weighted} & \textbf{0.1649} & \textbf{0.2307} \\
\texttt{fusion\_average} & 0.1610 & 0.2269 \\
\texttt{fusion\_concat} & 0.1519 & 0.2143 \\
\midrule
\multicolumn{3}{@{}l}{\textit{RGB frame-count control (fixed batch\,=\,32, 25 epochs)}} \\
\texttt{rgb\_n2\_ctrl} & 0.1316 & 0.1938 \\
\texttt{rgb\_n4\_ctrl} & 0.1549 & 0.2176 \\
\texttt{rgb\_n8\_ctrl} & 0.1620 & 0.2268 \\
\texttt{rgb\_n12\_ctrl} & 0.1672 & 0.2290 \\
\texttt{rgb\_n16\_ctrl} & 0.1690 & 0.2338 \\
\midrule
\multicolumn{3}{@{}l}{\textit{Dense 12\,fps sweeps (20 epochs)}} \\
\texttt{rgb\_n32\_dense\_e20} & \textbf{0.1701} & \textbf{0.2380} \\
\texttt{rgb\_n64\_dense\_e20} & 0.1600 & 0.2231 \\
\texttt{flow\_n32\_dense\_e20} & 0.1086 & 0.1707 \\
\texttt{flow\_n64\_dense\_e20} & 0.1069 & 0.1725 \\
\texttt{fusion\_average\_n32\_dense\_e20} & 0.1403 & 0.2052 \\
\texttt{fusion\_average\_n64\_dense\_e20} & 0.0964 & 0.1545 \\
\texttt{fusion\_concat\_n32\_dense\_e20} & 0.1119 & 0.1758 \\
\bottomrule
\end{tabular}
\end{table}

\subsection{VideoMAE RGB on Fir (completed)}

\begin{table}[h]
\centering
\caption{VideoMAE RGB results on Fir (protocol: \texttt{train\_test}).}
\label{tab:fir-videomae}
\begin{tabular}{@{}lrr@{}}
\toprule
Run & Macro mAP & Weighted mAP \\
\midrule
\textbf{\texttt{videomae\_vitb}} & \textbf{0.3259} & \textbf{0.4238} \\
\bottomrule
\end{tabular}
\end{table}

\noindent
At 16 frames per video, \texttt{videomae\_vitb} nearly doubles the best CNN result
(\texttt{rgb\_n32\_dense\_e20}, 0.1701 macro mAP) and exceeds the 16-frame RGB control
(\texttt{rgb\_n16\_ctrl}, 0.1690) by \textbf{+15.7 percentage points}.


\subsection{VideoMAE two-stream on Nibi (partial runs)}

\begin{table}[h]
\centering
\caption{VideoMAE two-stream results on Nibi (protocol: \texttt{train\_test}). Test macro and test weighted mAP from \texttt{per\_class\_ap.json}. All jobs terminated before 30 epochs; f64 evaluation did not complete.}
\label{tab:nibi-results}
\begin{tabular}{@{}lrrrr@{}}
\toprule
Run & Epochs & Batch & Test macro & Test wtd \\
\midrule
\texttt{videomae\_twostream\_vitb\_f16} & 8 / 30 & 6 & \textbf{0.2545} & \textbf{0.3481} \\
\texttt{videomae\_twostream\_vitb\_f32} & 4 / 30 & 2 & 0.0852 & 0.1406 \\
\texttt{videomae\_twostream\_vitb\_f64} & 3 / 30 & 1 &  0.0652  & 0.1091 \\
\bottomrule
\end{tabular}
\end{table}

%==============================================================================
\section{Discussion}
%==============================================================================

\subsection{CNN baselines}

\begin{itemize}[noitemsep]
  \item RGB substantially outperforms flow alone ($\sim$0.16 vs.\ $\sim$0.10 macro mAP).
  \item Among default-frame fusion strategies, \textbf{weighted} fusion is marginally best
        (0.1649), but gains over RGB-only (0.1616) are modest ($\sim$+0.3\,pp).
  \item Increasing RGB frames from 2 to 16 (control sweep) monotonically improves mAP
        (0.132 $\rightarrow$ 0.169), confirming temporal context matters.
  \item The best \textbf{completed} CNN result is \texttt{rgb\_n32\_dense\_e20} at
        \textbf{0.1701} macro mAP (12\,fps subset, 32 frames, 20 epochs).
  \item At 64 dense frames, RGB mAP \emph{decreases} slightly (0.1600), and fusion at
        64 frames collapses (0.0964), suggesting optimization difficulty or overfitting
        with very long clips under the CNN architecture.
\end{itemize}

\subsection{VideoMAE vs.\ CNN baselines}

\textbf{Completed RGB-only VideoMAE on Fir} (\texttt{videomae\_vitb}, 30 epochs,
16 frames) reaches \textbf{0.3259 macro mAP} and \textbf{0.4238 weighted mAP} ---
\textbf{+15.6 percentage points} over the best CNN (\texttt{rgb\_n32\_dense\_e20},
0.1701; +92\% relative) and \textbf{+15.7 pp} over \texttt{rgb\_n16\_ctrl} (0.1690)
at matched frame count. This confirms the gap reported to the supervisor
($\sim$32.6\% vs.\ $\sim$16.9\% at 16 frames).

Several factors likely contribute to this gap:

\begin{enumerate}[noitemsep]
  \item \textbf{Stronger pretrained representations:} ViT-Base trained on Kinetics-710 via
        masked video modeling vs.\ ImageNet-1K ResNet-50.
  \item \textbf{Native temporal modeling:} VideoMAE processes spatiotemporal tubelets
        jointly; CNNs pool frame-wise features with simple temporal averaging.
  \item \textbf{End-to-end clip encoding:} a single ViT consumes all 16 frames as one
        spatiotemporal volume rather than independent per-frame ResNet passes.
\end{enumerate}

\textbf{VideoMAE two-stream on Nibi} (\texttt{videomae\_twostream\_vitb\_f16}) reached
0.2545 macro mAP after only \textbf{8 of 30 epochs} --- above all CNN baselines but
\textbf{below} the completed RGB-only Fir run (0.3259). Two-stream training is harder
(two ViT backbones, OOM at longer clips) and was stopped early; a fair comparison
requires completing the two-stream recipe and an RGB-only baseline on Nibi at matched
settings.

From \texttt{history.json} on Nibi: f16 peaked at 0.2545 after 8 epochs (LR $10^{-3}$,
flow embed frozen 3 epochs); f32 reached 0.0852 after 4 epochs (batch 2, OOM); f64
reached 0.0652 after 3 epochs (batch 1).

\subsection{Interpreting the incomplete frame sweep on Nibi}

The apparent degradation from f16 $\rightarrow$ f32 $\rightarrow$ f64
(0.2545 $\rightarrow$ 0.0852 $\rightarrow$ 0.0652) \textbf{cannot be interpreted as a
frame-count trend} at this stage because:

\begin{itemize}[noitemsep]
  \item Training duration decreased (8, 4, 3 epochs), correlating with lower mAP.
  \item Batch sizes were reduced (6, 2, 1), slowing convergence and increasing gradient noise.
  \item Higher learning rate ($10^{-3}$) in the sweep may cause instability when the flow
        embed is unfrozen, especially at larger clip lengths.
\end{itemize}

A fair frame-count ablation requires \textbf{completed} 30-epoch runs with matched
optimization settings.

\subsection{Flow contribution}

On CNN baselines, adding flow via fusion yields only small gains over RGB-only at default
frame counts (weighted fusion: 0.1649 vs.\ RGB 0.1616). The completed
\texttt{videomae\_vitb} run establishes a strong RGB-only VideoMAE ceiling (0.3259);
once Nibi two-stream training finishes, comparing against
\texttt{jobs/nibi/train\_videomae\_rgb.sh} will show whether flow adds value at the
transformer level.

%==============================================================================
\section{Challenges Encountered}
%==============================================================================

\begin{enumerate}[leftmargin=*]

  \item \textbf{GPU memory (OOM) for long clips.}
  Two-stream VideoMAE loads \emph{two} ViT-Base backbones simultaneously. At 32 and 64
  frames, the default batch sizes (4 and 2) exceeded available H100 memory; batches were
  reduced to 2 and 1 respectively. CNN baselines mitigate this via temporal chunking
  (16 frames per forward) and gradient checkpointing; VideoMAE currently lacks an
  equivalent chunking strategy.

  \item \textbf{Job termination before target epochs.}
  All three Nibi sweep jobs stopped well short of 30 epochs (8, 4, 3). Likely causes
  include Slurm time limits (\texttt{--time=36:00:00} on the sweep array), node
  preemption, or training crashes. The f64 run failed before the post-training evaluation
  step, leaving no \texttt{per\_class\_ap.json}.

  \item \textbf{Learning-rate sensitivity.}
  The frame sweep uses peak LR $= 10^{-3}$. Internal documentation
  (\texttt{docs/nibi/REPRODUCE.md}) notes that $10^{-3}$ can produce NaN losses after
  unfreezing the flow patch embedding, whereas $5 \times 10^{-4}$ with gradient clipping
  is more stable. The sweep recipe and the reproducible recipe therefore differ; future
  runs should align with the stable recipe or use a lower LR on resume.
  
\end{enumerate}

%==============================================================================
\section{Next Steps}
%==============================================================================

\begin{enumerate}[leftmargin=*]
  \item \textbf{Resume and complete f16 training} to epoch~30 using
        \texttt{jobs/nibi/resume\_videomae\_twostream.sh} with a reduced learning rate
        ($10^{-4}$) to preserve the strong epoch-8 checkpoint.
  \item \textbf{Run the reproducible f16 recipe} from epoch~1
        (LR $= 5 \times 10^{-4}$, freeze flow embed 5 epochs, grad clip 1.0) as an
        independent confirmation.
  \item \textbf{Diagnose sweep job failures} via Slurm logs
        (\texttt{runs/nibi\_vmae2s\_frames-*.out}) and \texttt{sacct}.
  \item \textbf{Re-run f32/f64} with extended wall time, stable LR, and OOM-safe batch
        sizes after f16 is complete.
  
\end{enumerate}

\end{document}
